\section{Determining Cost and Budget}
\subsection{Estimate Activity Costs and the Total Budgeted Cost}
Every estimate is a retrospective estimate from commit history and team recall: hours from Table 4.1 multiplied by the rate card (Table 0.3). Materials, equipment, facilities, travel and subcontractors are ₹0, as the college provided the GPU and facility without charge and the data cost nothing. The GPU use is therefore uncharged rather than free of value. The reserve is held outside the TBC, so earned value is measured only against planned work.

\begin{longtable}{>{\raggedright\arraybackslash}p{1.01cm}>{\raggedright\arraybackslash}p{5.6cm}>{\raggedleft\arraybackslash}p{1.68cm}>{\raggedleft\arraybackslash}p{1.62cm}>{\raggedright\arraybackslash}p{4.71cm}>{\raggedleft\arraybackslash}p{1.68cm}}
\hline\rowcolor{navy}\thead{WBS number} & \thead{Work package} & \thead{Labour (₹)} & \thead{Materials / other (₹)} & \thead{Basis of estimate} & \thead{TBC (₹)} \\ \hline\endfirsthead
\hline\rowcolor{navy}\thead{WBS number} & \thead{Work package} & \thead{Labour (₹)} & \thead{Materials / other (₹)} & \thead{Basis of estimate} & \thead{TBC (₹)} \\ \hline\endhead
1.1 & Literature survey & ₹10,600 & ₹0 & 3 members for 2 weeks at about 8 h/week & ₹10,600 \\ \hline
1.2 & Requirements and architecture & ₹8,400 & ₹0 & Mostly PM; sized against the 15-day duration & ₹8,400 \\ \hline
2.1 & Platform scaffold and core & ₹6,200 & ₹0 & Size of the Phase 0-4 scaffold & ₹6,200 \\ \hline
2.2 & Backend API and database & ₹9,050 & ₹0 & API routers and DB schema count & ₹9,050 \\ \hline
2.3 & Dashboard v1 and redesign & ₹9,400 & ₹0 & Page count of the dashboard & ₹9,400 \\ \hline
3.1 & Cap dataset, detector and classifier iterations & ₹10,400 & ₹0 & Training runs expected, plus analysis & ₹10,400 \\ \hline
3.2 & Fill dataset and fill-level models & ₹5,600 & ₹0 & Two detectors and fill-ratio logic & ₹5,600 \\ \hline
3.3 & First integrated demo & ₹3,400 & ₹0 & Integration week & ₹3,400 \\ \hline
4.1 & Label: PatchCore baseline & ₹6,400 & ₹0 & One anomaly baseline & ₹6,400 \\ \hline
4.2 & Label: annotation, YOLO iteration, template matching & ₹9,000 & ₹0 & 528 crops plus localisation trials & ₹9,000 \\ \hline
4.3 & Label: EfficientAD, grayscale, calibration & ₹8,200 & ₹0 & 4 ROI models with calibration & ₹8,200 \\ \hline
4.4 & Label dashboard & ₹5,000 & ₹0 & Single-purpose dashboard & ₹5,000 \\ \hline
5.1 & Cap teacher-student distillation & ₹7,050 & ₹0 & Teacher and student training watch & ₹7,050 \\ \hline
5.2 & Station 1 production pipeline & ₹10,050 & ₹0 & Six pipeline layers & ₹10,050 \\ \hline
5.3 & Docker and TensorRT deployment & ₹7,250 & ₹0 & Compose, engine compile, stress test & ₹7,250 \\ \hline
5.4 & Audit, Replay, Active Learning workstations & ₹7,250 & ₹0 & Three operator pages & ₹7,250 \\ \hline
5.5 & Station 2 integration, env-driven paths & ₹5,250 & ₹0 & Mirror of Station 1 layout & ₹5,250 \\ \hline
5.6 & GPU verification and handover & ₹3,400 & ₹0 & Checks on target PC and handover & ₹3,400 \\ \hline
6.1 & Thesis, IEEE paper, presentation & ₹10,600 & ₹0 & Thesis, paper and slides, two weeks & ₹10,600 \\ \hline
7.1 & Weekly guide reviews & ₹32,000 & ₹0 & 32 h planned at ₹1,000/h & ₹32,000 \\ \hline
\textbf{} & \textbf{Total budgeted cost (TBC)} & \textbf{₹1,74,500} & \textbf{₹0} & \textbf{Sum of work packages} & \textbf{₹1,74,500} \\ \hline
\textbf{} & \textbf{Contingency reserve (held separately)} & \textbf{} & \textbf{} & \textbf{10\% of TBC} & \textbf{₹17,450} \\ \hline
\textbf{} & \textbf{Funding limit} & \textbf{} & \textbf{} & \textbf{TBC + reserve} & \textbf{₹1,91,950} \\ \hline
\end{longtable}\nopagebreak
\vspace{-1.2em}\nopagebreak\tcap{Table 5.1 - Activity cost estimates and total budgeted cost (all amounts in ₹)}


\subsection{Time-Phased Budget (Cumulative Budgeted Cost)}
Each package is spread evenly over its planned working days. Rows use activity letters from Table 4.1. Amounts are in ₹.

{\footnotesize
\begin{longtable}{>{\raggedright\arraybackslash}p{1.04cm}>{\raggedleft\arraybackslash}p{1.41cm}>{\raggedleft\arraybackslash}p{1.41cm}>{\raggedleft\arraybackslash}p{1.41cm}>{\raggedleft\arraybackslash}p{1.41cm}>{\raggedleft\arraybackslash}p{1.41cm}>{\raggedleft\arraybackslash}p{1.41cm}>{\raggedleft\arraybackslash}p{1.41cm}>{\raggedleft\arraybackslash}p{1.41cm}>{\raggedleft\arraybackslash}p{1.41cm}>{\raggedleft\arraybackslash}p{1.57cm}}
\hline\rowcolor{navy}\tablelist{Table 5.2 - Time-phased budget}\thead{WP} & \thead{P1 Feb} & \thead{P2 Mar} & \thead{P3 Apr} & \thead{P4 May} & \thead{P5 Jun} & \thead{P6 Jul} & \thead{P7 Aug} & \thead{P8 Sep} & \thead{P9 Oct} & \thead{Total (₹)} \\ \hline\endfirsthead
\hline\rowcolor{navy}\thead{WP} & \thead{P1 Feb} & \thead{P2 Mar} & \thead{P3 Apr} & \thead{P4 May} & \thead{P5 Jun} & \thead{P6 Jul} & \thead{P7 Aug} & \thead{P8 Sep} & \thead{P9 Oct} & \thead{Total (₹)} \\ \hline\endhead
A & 10,600 &  &  &  &  &  &  &  &  & 10,600 \\ \hline
B & 5,600 & 2,800 &  &  &  &  &  &  &  & 8,400 \\ \hline
C &  & 6,200 &  &  &  &  &  &  &  & 6,200 \\ \hline
D &  & 3,168 & 5,882 &  &  &  &  &  &  & 9,050 \\ \hline
E &  & 2,193 & 6,894 & 313 &  &  &  &  &  & 9,400 \\ \hline
F &  & 7,072 & 3,328 &  &  &  &  &  &  & 10,400 \\ \hline
G &  &  & 5,600 &  &  &  &  &  &  & 5,600 \\ \hline
H &  &  &  & 3,400 &  &  &  &  &  & 3,400 \\ \hline
I &  &  &  & 10,600 &  &  &  &  &  & 10,600 \\ \hline
J &  &  &  & 1,280 & 5,120 &  &  &  &  & 6,400 \\ \hline
K &  &  &  &  & 1,200 & 7,800 &  &  &  & 9,000 \\ \hline
L &  &  &  &  &  & 5,467 & 2,733 &  &  & 8,200 \\ \hline
M &  &  &  &  &  &  & 5,000 &  &  & 5,000 \\ \hline
N &  &  &  &  &  &  & 7,050 &  &  & 7,050 \\ \hline
O &  &  &  &  &  &  & 6,030 & 4,020 &  & 10,050 \\ \hline
P &  &  &  &  &  &  &  & 7,250 &  & 7,250 \\ \hline
Q &  &  &  &  &  &  &  & 7,250 &  & 7,250 \\ \hline
R &  &  &  &  &  &  &  & 5,250 &  & 5,250 \\ \hline
S &  &  &  &  &  &  &  & 3,400 &  & 3,400 \\ \hline
7.1 & 3,699 & 4,070 & 4,069 & 3,885 & 4,069 & 4,254 & 3,885 & 4,069 &  & 32,000 \\ \hline
\textbf{Period} & \textbf{19,899} & \textbf{25,503} & \textbf{25,773} & \textbf{19,478} & \textbf{10,389} & \textbf{17,521} & \textbf{24,698} & \textbf{31,239} & \textbf{0} & \textbf{174,500} \\ \hline
\textbf{CBC} & \textbf{19,899} & \textbf{45,402} & \textbf{71,175} & \textbf{90,653} & \textbf{101,042} & \textbf{118,563} & \textbf{143,261} & \textbf{174,500} & \textbf{174,500} & \textbf{} \\ \hline
\end{longtable}\nopagebreak
\vspace{-1.2em}\nopagebreak\tcap{Table 5.2 - Time-phased budget; the last row is the cumulative budgeted cost (CBC)}
}


\subsection{Actual Cost and Committed Cost}
Spent means hours from the reconstruction valued at the rate card. No cash was paid and nothing was ordered, so committed cost is ₹0 throughout.

\begin{longtable}{>{\raggedright\arraybackslash}p{1.01cm}>{\raggedright\arraybackslash}p{7.87cm}>{\raggedleft\arraybackslash}p{1.91cm}>{\raggedleft\arraybackslash}p{1.91cm}>{\raggedleft\arraybackslash}p{1.69cm}>{\raggedleft\arraybackslash}p{1.91cm}}
\hline\rowcolor{navy}\tablelist{Table 5.3 - Actual and committed cost as of the Status Date}\thead{WBS} & \thead{Work package} & \thead{CBC to date (₹)} & \thead{Spent (₹, notional)} & \thead{Committed (₹)} & \thead{CAC (₹)} \\ \hline\endfirsthead
\hline\rowcolor{navy}\thead{WBS} & \thead{Work package} & \thead{CBC to date (₹)} & \thead{Spent (₹, notional)} & \thead{Committed (₹)} & \thead{CAC (₹)} \\ \hline\endhead
1.1 & Literature survey & ₹10,600 & ₹10,600 & ₹0 & ₹10,600 \\ \hline
1.2 & Requirements and architecture & ₹8,400 & ₹7,900 & ₹0 & ₹7,900 \\ \hline
2.1 & Platform scaffold and core & ₹6,200 & ₹6,200 & ₹0 & ₹6,200 \\ \hline
2.2 & Backend API and database & ₹9,050 & ₹8,850 & ₹0 & ₹8,850 \\ \hline
2.3 & Dashboard v1 and redesign & ₹9,400 & ₹11,000 & ₹0 & ₹11,000 \\ \hline
3.1 & Cap dataset, detector and classifier iterations & ₹10,400 & ₹12,000 & ₹0 & ₹12,000 \\ \hline
3.2 & Fill dataset and fill-level models & ₹5,600 & ₹6,600 & ₹0 & ₹6,600 \\ \hline
3.3 & First integrated demo & ₹3,400 & ₹3,400 & ₹0 & ₹3,400 \\ \hline
4.1 & Label: PatchCore baseline & ₹6,400 & ₹8,000 & ₹0 & ₹8,000 \\ \hline
4.2 & Label: annotation, YOLO iteration, template matching & ₹9,000 & ₹11,600 & ₹0 & ₹11,600 \\ \hline
4.3 & Label: EfficientAD, grayscale, calibration & ₹8,200 & ₹8,200 & ₹0 & ₹8,200 \\ \hline
4.4 & Label dashboard & ₹5,000 & ₹6,800 & ₹0 & ₹6,800 \\ \hline
5.1 & Cap teacher-student distillation & ₹7,050 & ₹7,550 & ₹0 & ₹7,550 \\ \hline
5.2 & Station 1 production pipeline & ₹10,050 & ₹10,050 & ₹0 & ₹10,050 \\ \hline
5.3 & Docker and TensorRT deployment & ₹7,250 & ₹8,750 & ₹0 & ₹8,750 \\ \hline
5.4 & Audit, Replay, Active Learning workstations & ₹7,250 & ₹7,850 & ₹0 & ₹7,850 \\ \hline
5.5 & Station 2 integration, env-driven paths & ₹5,250 & ₹5,250 & ₹0 & ₹5,250 \\ \hline
5.6 & GPU verification and handover & ₹3,400 & ₹4,200 & ₹0 & ₹4,200 \\ \hline
6.1 & Thesis, IEEE paper, presentation & ₹10,600 & ₹9,700 & ₹0 & ₹9,700 \\ \hline
7.1 & Weekly guide reviews & ₹32,000 & ₹35,000 & ₹0 & ₹35,000 \\ \hline
\textbf{} & \textbf{Total} & \textbf{₹1,74,500} & \textbf{₹1,89,500} & \textbf{₹0} & \textbf{₹1,89,500} \\ \hline
\end{longtable}\nopagebreak
\vspace{-1.2em}\nopagebreak\tcap{Table 5.3 - Actual and committed cost as of the Status Date}


\subsection{Earned Value of the Work Performed}
Percent complete at the Status Date was judged by Sushanth S as PM, against an artifact that can be checked in the repository, not by opinion. For monthly data, an in-progress package is credited with the elapsed share of its actual window, rounded to the nearest 5\% and held below 100\% until its finishing commit.

\begin{longtable}{>{\raggedright\arraybackslash}p{1.0cm}>{\raggedright\arraybackslash}p{6.21cm}>{\raggedleft\arraybackslash}p{1.66cm}>{\centering\arraybackslash}p{1.44cm}>{\raggedright\arraybackslash}p{4.32cm}>{\raggedleft\arraybackslash}p{1.67cm}}
\hline\rowcolor{navy}\tablelist{Table 5.4 - Cumulative earned value}\thead{WBS} & \thead{Work package} & \thead{TBC (₹)} & \thead{Percent complete} & \thead{How judged} & \thead{CEV (₹)} \\ \hline\endfirsthead
\hline\rowcolor{navy}\thead{WBS} & \thead{Work package} & \thead{TBC (₹)} & \thead{Percent complete} & \thead{How judged} & \thead{CEV (₹)} \\ \hline\endhead
1.1 & Literature survey & ₹10,600 & 100\% & Reference list in the thesis & ₹10,600 \\ \hline
1.2 & Requirements and architecture & ₹8,400 & 100\% & Planning documents exist & ₹8,400 \\ \hline
2.1 & Platform scaffold and core & ₹6,200 & 100\% & Code in repo & ₹6,200 \\ \hline
2.2 & Backend API and database & ₹9,050 & 100\% & API and schema commits & ₹9,050 \\ \hline
2.3 & Dashboard v1 and redesign & ₹9,400 & 100\% & Dashboard commits & ₹9,400 \\ \hline
3.1 & Cap dataset, detector and classifier iterations & ₹10,400 & 100\% & Models in repo, metrics logged & ₹10,400 \\ \hline
3.2 & Fill dataset and fill-level models & ₹5,600 & 100\% & Fill models run in pipeline & ₹5,600 \\ \hline
3.3 & First integrated demo & ₹3,400 & 100\% & Commit 8c54630 & ₹3,400 \\ \hline
4.1 & Label: PatchCore baseline & ₹6,400 & 100\% & PatchCore scripts, membank & ₹6,400 \\ \hline
4.2 & Label: annotation, YOLO iteration, template matching & ₹9,000 & 100\% & Crops, localisation committed & ₹9,000 \\ \hline
4.3 & Label: EfficientAD, grayscale, calibration & ₹8,200 & 100\% & Thresholds, crop test records & ₹8,200 \\ \hline
4.4 & Label dashboard & ₹5,000 & 100\% & Separate dashboard runs & ₹5,000 \\ \hline
5.1 & Cap teacher-student distillation & ₹7,050 & 100\% & results.csv, 126 epochs & ₹7,050 \\ \hline
5.2 & Station 1 production pipeline & ₹10,050 & 100\% & test\_step1-6 scripts & ₹10,050 \\ \hline
5.3 & Docker and TensorRT deployment & ₹7,250 & 100\% & Compose files, stress test & ₹7,250 \\ \hline
5.4 & Audit, Replay, Active Learning workstations & ₹7,250 & 100\% & Commit c8f0524 & ₹7,250 \\ \hline
5.5 & Station 2 integration, env-driven paths & ₹5,250 & 100\% & test\_system\_e2e.py & ₹5,250 \\ \hline
5.6 & GPU verification and handover & ₹3,400 & 100\% & Team demonstration & ₹3,400 \\ \hline
6.1 & Thesis, IEEE paper, presentation & ₹10,600 & 100\% & Thesis and plagiarism report & ₹10,600 \\ \hline
7.1 & Weekly guide reviews & ₹32,000 & 100\% & Level of effort & ₹32,000 \\ \hline
\textbf{} & \textbf{Total} & \textbf{₹1,74,500} & \textbf{100\%} & \textbf{} & \textbf{₹1,74,500} \\ \hline
\end{longtable}\nopagebreak
\vspace{-1.2em}\nopagebreak\tcap{Table 5.4 - Cumulative earned value}


\subsection{Analyse Cost Performance}
\simpletable{L{1.8cm}L{4.2cm}L{6.0cm}L{4.7cm}}{\thead{Measure} & \thead{Formula} & \thead{Value} & \thead{What it tells us}}{
\textbf{TBC} & Sum of work package budgets & \vTBC{} & Performance baseline \\ \hline
\textbf{CBC} & Budgeted for work scheduled to date & \vCBC{} & All work was due by 30 Sep \\ \hline
\textbf{CAC} & Spent plus committed & \vCAC{} + ₹0 = \vCAC{} & Notional effort used \\ \hline
\textbf{CEV} & TBC $\times$ percent complete & \vTBC{} $\times$ 100\% = \vCEV{} & All planned work earned \\ \hline
\textbf{CPI} & CEV / CAC & \vCEV{} / \vCAC{} = \vCPI{} & ₹\vCPI{} of value per ₹1 of effort \\ \hline
\textbf{CV} & CEV - CAC & \vCEV{} - \vCAC{} = \vCV{} & Over budget, within reserve \\ \hline
}{Table 5.5 - Cost performance as of the Status Date}

\begin{center}\resizebox{\textwidth}{!}{\begin{tikzpicture}[font=\footnotesize]
\draw[gridline] (0,0.00) -- (13.20,0.00);
\node[anchor=east,text=inksecondary] at (-0.05,0.00) {0.0};
\draw[gridline] (0,1.75) -- (13.20,1.75);
\node[anchor=east,text=inksecondary] at (-0.05,1.75) {0.5};
\draw[gridline] (0,3.50) -- (13.20,3.50);
\node[anchor=east,text=inksecondary] at (-0.05,3.50) {1.0};
\draw[gridline] (0,5.25) -- (13.20,5.25);
\node[anchor=east,text=inksecondary] at (-0.05,5.25) {1.5};
\draw[gridline] (0,7.00) -- (13.20,7.00);
\node[anchor=east,text=inksecondary] at (-0.05,7.00) {2.0};
\node[rotate=90,inner sep=0pt,text=inksecondary] at (-0.95,3.50) {₹ lakh, cumulative};
\node[anchor=north,text=inksecondary] at (1.44,-0.05) {Feb};
\node[anchor=north,text=inksecondary] at (2.89,-0.05) {Mar};
\node[anchor=north,text=inksecondary] at (4.33,-0.05) {Apr};
\node[anchor=north,text=inksecondary] at (5.78,-0.05) {May};
\node[anchor=north,text=inksecondary] at (7.22,-0.05) {Jun};
\node[anchor=north,text=inksecondary] at (8.67,-0.05) {Jul};
\node[anchor=north,text=inksecondary] at (10.11,-0.05) {Aug};
\node[anchor=north,text=inksecondary] at (11.56,-0.05) {Sep};
\node[anchor=north,text=inksecondary] at (13.00,-0.05) {Oct};
\draw[inksecondary] (0,0) -- (13.20,0);
\draw[marker] (8.67,0) -- (8.67,6.30);
\node[anchor=west,text=inksecondary] at (8.67,0.45) {31 Jul checkpoint};
\draw[curveCBC] (1.44,0.70) -- (2.89,1.59) -- (4.33,2.49) -- (5.78,3.17) -- (7.22,3.54) -- (8.67,4.15) -- (10.11,5.01) -- (11.56,6.11) -- (13.00,6.11);
\node[curveCBCmark] at (1.44,0.70) {};
\node[curveCBCmark] at (2.89,1.59) {};
\node[curveCBCmark] at (4.33,2.49) {};
\node[curveCBCmark] at (5.78,3.17) {};
\node[curveCBCmark] at (7.22,3.54) {};
\node[curveCBCmark] at (8.67,4.15) {};
\node[curveCBCmark] at (10.11,5.01) {};
\node[curveCBCmark] at (11.56,6.11) {};
\node[curveCBCmark] at (13.00,6.11) {};
\draw[curveCAC] (1.44,0.68) -- (2.89,1.43) -- (4.33,2.45) -- (5.78,3.21) -- (7.22,3.52) -- (8.67,4.33) -- (10.11,4.88) -- (11.56,5.83) -- (13.00,6.63);
\node[curveCACmark] at (1.44,0.68) {};
\node[curveCACmark] at (2.89,1.43) {};
\node[curveCACmark] at (4.33,2.45) {};
\node[curveCACmark] at (5.78,3.21) {};
\node[curveCACmark] at (7.22,3.52) {};
\node[curveCACmark] at (8.67,4.33) {};
\node[curveCACmark] at (10.11,4.88) {};
\node[curveCACmark] at (11.56,5.83) {};
\node[curveCACmark] at (13.00,6.63) {};
\draw[curveCEV] (1.44,0.68) -- (2.89,1.41) -- (4.33,2.32) -- (5.78,3.08) -- (7.22,3.34) -- (8.67,4.03) -- (10.11,4.52) -- (11.56,5.39) -- (13.00,6.11);
\node[curveCEVmark] at (1.44,0.68) {};
\node[curveCEVmark] at (2.89,1.41) {};
\node[curveCEVmark] at (4.33,2.32) {};
\node[curveCEVmark] at (5.78,3.08) {};
\node[curveCEVmark] at (7.22,3.34) {};
\node[curveCEVmark] at (8.67,4.03) {};
\node[curveCEVmark] at (10.11,4.52) {};
\node[curveCEVmark] at (11.56,5.39) {};
\node[curveCEVmark] at (13.00,6.11) {};
\draw[curveCBC] (0.25,6.65) -- (0.95,6.65);
\node[curveCBCmark] at (0.60,6.65) {};
\node[anchor=west] at (1.00,6.65) {CBC (budgeted)};
\draw[curveCAC] (3.75,6.65) -- (4.45,6.65);
\node[curveCACmark] at (4.10,6.65) {};
\node[anchor=west] at (4.50,6.65) {CAC (actual)};
\draw[curveCEV] (7.25,6.65) -- (7.95,6.65);
\node[curveCEVmark] at (7.60,6.65) {};
\node[anchor=west] at (8.00,6.65) {CEV (earned)};
\node[anchor=west,inner sep=0pt] at (13.15,6.78) {CAC};
\node[anchor=west,inner sep=0pt] at (13.15,5.96) {CEV, CBC};
\end{tikzpicture}
}\end{center}
\tcap{Figure 5.1 - Cumulative budgeted, actual and earned value curves. Table 5.8 gives the values.}

\textbf{What the picture shows.} CAC is above CEV from March onward and the gap widens, so we were consistently paying slightly more for each unit of work. CEV falls well behind CBC during August and September, which is when the schedule slip appears as money: planned work was not being earned. In October CBC is flat because all work was due by then, and CEV catches up with it at 100\%, but CAC does not come back down.

\subsection{Forecast the Cost at Completion}
At the Status Date the work is done, so both forecasts equal the final cost and the TCPI has no remaining work to apply to.

\simpletable{L{3.32cm}L{4.33cm}L{4.02cm}L{5.03cm}}{\thead{Measure} & \thead{Formula} & \thead{Value} & \thead{Interpretation}}{
\textbf{FCAC, method 1} & TBC / CPI & \vTBC{} / \vCPI{} = \vFCACone{} & Equals CAC at completion \\ \hline
\textbf{FCAC, method 2} & CAC + (TBC - CEV) & \vCAC{} + ₹0 = \vFCACtwo{} & Equals CAC at completion \\ \hline
\textbf{Overrun or underrun} & TBC - FCAC & \vTBC{} - \vCAC{} = \vOverrun{} & Overrun; \vReserveUsedPct{}\% of reserve used \\ \hline
\textbf{TCPI} & (TBC - CEV) / (TBC - CAC) & ₹0 / \vOverrun{} = 0 & Finishing within TBC is no longer possible \\ \hline
}{Table 5.6 - Forecasted cost at completion}

\textbf{Interim forecast, 31 July 2026.} A completed project cannot show a live forecast, so the table below applies the same methods at the last month before the deployment slips, when \viPct{}\% of the value was earned.

\simpletable{L{4.23cm}L{5.03cm}L{3.02cm}L{4.42cm}}{\thead{Measure at 31 July} & \thead{Working} & \thead{Value} & \thead{Reading}}{
CBC, CAC, CEV & From Table 5.8 & \viCBC{}, \viCAC{}, \viCEV{} & Slightly behind and over \\ \hline
CPI & \viCEV{} / \viCAC{} & \viCPI{} & Below 1 \\ \hline
CV & \viCEV{} - \viCAC{} & \viCV{} & Over budget \\ \hline
FCAC, method 1 & \vTBC{} / \viCPI{} & \viFCACone{} & Rest at today's efficiency \\ \hline
FCAC, method 2 & \viCAC{} + (\vTBC{} - \viCEV{}) & \viFCACtwo{} & Rest exactly to budget \\ \hline
TCPI & (\vTBC{} - \viCEV{}) / (\vTBC{} - \viCAC{}) & \viTCPI{} & Above our CPI, so unlikely \\ \hline
}{Table 5.7 - Interim cost forecast, 31 July 2026}

Method 1 was the more realistic assumption, because the overruns came from recurring causes, model iteration and deployment rework, not one-off events. The final cost of \vCAC{} was in fact above both forecasts, as the deployment rework was still to come. A TCPI above 1 meant that staying within the TBC would have needed better efficiency than we had achieved, which was unlikely, so the reserve was the realistic cover.

\textbf{If CPI had fallen to 0.80:} FCAC = \vTBC{} / 0.80 = \vFCACeighty{}, an overrun of \vOverrunEighty{}, more than the \vRESERVE{} reserve. We would have told the PM to cut remaining scope, since extra funding has no meaning here and each extra rupee is really student hours.

\begin{longtable}{>{\raggedright\arraybackslash}p{1.53cm}>{\raggedleft\arraybackslash}p{2.68cm}>{\raggedleft\arraybackslash}p{2.68cm}>{\raggedleft\arraybackslash}p{2.68cm}>{\raggedleft\arraybackslash}p{1.66cm}>{\raggedleft\arraybackslash}p{2.43cm}>{\raggedleft\arraybackslash}p{2.44cm}}
\hline\rowcolor{navy}\thead{Month end} & \thead{CBC} & \thead{CAC} & \thead{CEV} & \thead{CPI} & \thead{CV} & \thead{CEV - CBC} \\ \hline\endfirsthead
\hline\rowcolor{navy}\thead{Month end} & \thead{CBC} & \thead{CAC} & \thead{CEV} & \thead{CPI} & \thead{CV} & \thead{CEV - CBC} \\ \hline\endhead
Feb & ₹19,899 & ₹19,344 & ₹19,411 & 1.003 & ₹67 & -₹488 \\ \hline
Mar & ₹45,402 & ₹40,778 & ₹40,354 & 0.990 & -₹424 & -₹5,048 \\ \hline
Apr & ₹71,175 & ₹70,138 & ₹66,303 & 0.945 & -₹3,835 & -₹4,872 \\ \hline
May & ₹90,653 & ₹91,826 & ₹87,891 & 0.957 & -₹3,935 & -₹2,762 \\ \hline
Jun & ₹1,01,042 & ₹1,00,657 & ₹95,417 & 0.948 & -₹5,240 & -₹5,625 \\ \hline
Jul & ₹1,18,563 & ₹1,23,774 & ₹1,15,260 & 0.931 & -₹8,514 & -₹3,303 \\ \hline
Aug & ₹1,43,261 & ₹1,39,465 & ₹1,29,176 & 0.926 & -₹10,289 & -₹14,085 \\ \hline
Sep & ₹1,74,500 & ₹1,66,519 & ₹1,54,134 & 0.926 & -₹12,385 & -₹20,366 \\ \hline
Oct & ₹1,74,500 & ₹1,89,500 & ₹1,74,500 & 0.921 & -₹15,000 & ₹0 \\ \hline
\end{longtable}\nopagebreak
\vspace{-1.2em}\nopagebreak\tcap{Table 5.8 - Monthly control data behind Figure 5.1 (October is the Status Date)}

